% deskew_timeline.tex -- LiDAR-inertial de-skew timeline (FAST-LIO2 style).
% Time axis t_k..t_{k+1}; dense IMU ticks; LiDAR points at their own timestamps;
% state propagation; backward-propagation to the scan-end frame; scan-period brace.
% Has Chinese label text -> needs ctex + xelatex.
\documentclass[border=8pt]{standalone}
\usepackage{ctex}
\usepackage{tikz}
\usepackage{amsmath}
\usetikzlibrary{arrows.meta, decorations.pathreplacing, positioning, calc}
\begin{document}
\begin{tikzpicture}[
    >=Stealth, line cap=round, font=\small,
    imu/.style={blue!70!black, thick},
    lidar/.style={red!80!black},
    prop/.style={green!50!black, -{Stealth[length=2mm]}, thick}]

  \def\xk{0}\def\xkp{11}
  % ---- time axis ----
  \draw[->, thick] (\xk-0.6,0) -- (\xkp+1,0) node[right]{$t$};
  \foreach \x/\lab in {\xk/{t_k}, \xkp/{t_{k+1}}}
      \draw[thick] (\x,0.12)--(\x,-0.12) node[below]{$\lab$};

  % ---- IMU samples: dense ticks (high rate) ----
  \foreach \x in {0.7,1.7,2.7,3.7,4.7,5.7,6.7,7.7,8.7,9.7,10.5}
      \draw[imu] (\x,0)--(\x,0.45);
  \node[imu] at (5.5,0.85) {IMU $(\boldsymbol{\omega}_i,\mathbf a_i)@\tau_i$ (高频)};

  % ---- state propagation (top lane) ----
  \def\ys{3.4}
  \node[left] at (\xk,\ys){$\mathbf x_k$};  \node[right] at (\xkp,\ys){$\mathbf x_{k+1}$};
  \draw[prop] (\xk,\ys) .. controls (4,\ys+0.5) and (7,\ys-0.3) .. (\xkp,\ys);
  \node[green!45!black, above] at (5.5,\ys+0.35){IMU 积分传播状态};

  % ---- LiDAR points at own timestamps (middle lane) ----
  \def\yl{1.8}
  \node[lidar, left] at (\xk-0.3,\yl){LiDAR};
  \foreach \x/\j in {0.9/1, 2.5/2, 4.3/3, 6.1/4, 7.9/5, 9.6/6}{
      \fill[lidar] (\x,\yl) circle (2pt);
      \node[lidar, font=\scriptsize, above] at (\x,\yl+0.04){$\mathbf p_{\j}$};
      \draw[lidar, dotted] (\x,\yl)--(\x,0);}     % drop to its timestamp rho_j

  % ---- backward propagation: pull each point to the scan-end frame ----
  \draw[dashed, thick] (\xkp,-0.3)--(\xkp,\ys+0.2);  % de-skew reference @ t_{k+1}
  \foreach \x in {0.9,2.5,4.3,6.1,7.9}
      \draw[prop, densely dashed] (\x,\yl) .. controls (\x+1.6,\yl+0.6) .. (\xkp,\yl);
  \node[green!45!black, align=center, font=\scriptsize]
      at (8.7,\yl+1.0){反向传播 $\mathbf T_j$\\ 投到 $t_{k+1}$ 帧};

  % ---- brace: one scan period ----
  \draw[decorate, decoration={brace, amplitude=6pt, mirror}, thick]
      (\xk,-0.7)--(\xkp,-0.7) node[midway, below=8pt]{一帧扫描周期（去畸变区间）};
\end{tikzpicture}
\end{document}
